\documentclass[10pt,twocolumn]{article}
\usepackage[margin=1.6cm,top=1.8cm,bottom=1.8cm]{geometry}
\usepackage{amsmath,amssymb,amsthm}
\usepackage{booktabs}
\usepackage{microtype}
\usepackage{xcolor}
\usepackage{enumitem}
\usepackage{titlesec}

% ---------- colours ----------
\definecolor{lblue}{RGB}{55,138,221}
\definecolor{rred}{RGB}{226,75,74}

% ---------- layout ----------
\titleformat{\section}{\normalfont\bfseries\small}{}{0em}{}
\titlespacing{\section}{0pt}{5pt}{2pt}
\setlength{\columnsep}{14pt}
\setlength{\parskip}{2pt}
\setlength{\parindent}{8pt}

% ---------- theorem environments ----------
\theoremstyle{plain}
\newtheorem{theorem}{Theorem}
\newtheorem{proposition}[theorem]{Proposition}
\newtheorem{corollary}[theorem]{Corollary}
\theoremstyle{definition}
\newtheorem{definition}[theorem]{Definition}
\theoremstyle{remark}
\newtheorem{remark}[theorem]{Remark}

% ---------- macros ----------
\newcommand{\Lopt}[1]{{\color{lblue}#1}}
\newcommand{\Ropt}[1]{{\color{rred}#1}}
\newcommand{\val}[1]{\mathcal{G}(#1)}
\newcommand{\temp}[1]{t(#1)}
\newcommand{\mean}[1]{m(#1)}
\newcommand{\edgeL}{\overset{\Lopt{L}}{\longrightarrow}}
\newcommand{\edgeR}{\overset{\Ropt{R}}{\longrightarrow}}

% ---------- manual title block ----------
\newcommand{\papertitle}{%
  \begingroup
  \centering
  {\large\bfseries Partisan Tree Ascent: Edge-Labelled Games
    and Their Combinatorial Game Values}\par
  \vspace{4pt}
  {\small A.\ Researcher\,$^\dagger$\quad B.\ Co-author\,$^\dagger$\quad
    \textit{$^\dagger$ Department of Mathematics}}\par
  \vspace{6pt}
  \endgroup
}

\begin{document}

\twocolumn[{%
  \papertitle
  \noindent\small\itshape
  We define \emph{Partisan Tree Ascent}, a two-player perfect-information
  game played on finite rooted binary trees whose edges carry labels
  $\Lopt{L}$ or $\Ropt{R}$.  Left (resp.\ Right) moves the shared token
  along any adjacent $\Lopt{L}$-edge (resp.\ $\Ropt{R}$-edge), always away
  from the root; the player unable to move loses (normal play convention).
  We compute the CGT values of representative trees under four
  edge-labelling regimes, prove a monotonicity theorem for pure-colour
  paths, and derive a temperature formula that exactly determines
  first-move urgency in compound positions.
  \par\vspace{6pt}
}]

% ============================================================
\section{1.\enspace Game definition}
% ============================================================

\begin{definition}
A \emph{Partisan Tree Ascent position} is a pair $(T,v)$ where $T$ is a
finite rooted binary tree with each edge labelled $\Lopt{L}$ or $\Ropt{R}$,
and $v\in V(T)$ is the current token location.
\textbf{Left} moves $v$ to any child reachable via an $\Lopt{L}$-edge;
\textbf{Right} moves via any $\Ropt{R}$-edge.
The player without a legal move \textbf{loses}.
\end{definition}

The recursive CGT value at node $v$ is
\[
  \val{v} =
  \bigl\{\val{c}\mid c\in\Lopt{\mathcal{L}}(v)
  \;\big|\;
  \val{c}\mid c\in\Ropt{\mathcal{R}}(v)\bigr\},
\]
with $\val{\ell}=\{|\}=0$ at every leaf $\ell$.

% ============================================================
\section{2.\enspace Values by edge-labelling type}
% ============================================================

Table~\ref{tab:values} catalogues values for representative paths and
branching nodes.

\begin{table}[h]
\centering
\small
\setlength{\tabcolsep}{4pt}
\caption{CGT values of selected positions. Temperature $t=-\infty$
  denotes a number (cold game).}
\label{tab:values}
\vspace{2pt}
\begin{tabular}{@{}llcc@{}}
\toprule
Position & Type & $G$ & $t(G)$ \\
\midrule
$\Lopt{L}^1$                 & I   & $1$             & $-\infty$ \\
$\Ropt{R}^1$                 & I   & $-1$            & $-\infty$ \\
$\Lopt{L}^2$                 & I   & $2$             & $-\infty$ \\
$\Lopt{L}\Ropt{R}$           & II  & $0$             & $-\infty$ \\
$\Lopt{L}\Ropt{R}\Lopt{L}$   & II  & $1$             & $-\infty$ \\
$\{0\mid 0\}$                & III & $*$             & $0$       \\
$\{0\mid 1\}$                & III & $\tfrac{1}{2}$  & $-\infty$ \\
$\{1\mid{-1}\}$              & III & $\pm 1$         & $1$       \\
$\{2\mid{-2}\}$              & III & $\pm 2$         & $2$       \\
$\{0\mid *\}$                & IV  & $\uparrow$      & $0$       \\
$\{*\mid 0\}$                & IV  & $\downarrow$    & $0$       \\
$\{\uparrow\mid *\}$         & IV  & $\Uparrow$      & $0$       \\
\bottomrule
\end{tabular}
\end{table}

% ============================================================
\section{3.\enspace A monotonicity theorem}
% ============================================================

\begin{theorem}[Monochromatic path values]\label{thm:mono}
Let $P^n_L$ denote the path $v_0\edgeL v_1\edgeL\cdots\edgeL v_n$ rooted
at $v_0$.  Then $\val{P^n_L}=n$ for all $n\ge 0$.
Symmetrically, $\val{P^n_R}=-n$.
\end{theorem}

\begin{proof}
By induction on $n$.  Base: $\val{P^0_L}=\{|\}=0$.
Inductive step: Left's unique option from $v_0$ leads to the tail
$P^{n-1}_L$ with value $n-1$ by hypothesis; Right has no options.
Hence $\val{P^n_L}=\{n-1\mid\,\}$.
By the \emph{simplicity rule}, $\{n-1\mid\,\}=n$, since $n$ is the
simplest number strictly greater than $n-1$ with no Right upper bound.
\end{proof}

\begin{corollary}\label{cor:outcome}
$P^n_L$ is a Left win unconditionally for all $n\ge 1$,
regardless of who moves first.
\end{corollary}

\begin{proof}
For $n\ge 1$, $\val{P^n_L}=n>0$, so Left wins under optimal play.
Since Right cannot move at all, the outcome class is $\mathcal{L}$
irrespective of turn order---a stronger claim than ``first-player wins.''
\end{proof}

\begin{remark}
Purely monochromatic positions contribute only a constant integer to
$\val{G_1+\cdots+G_k}$ and offer no strategic choice.  All
game-theoretic content lies in branching and mixed-colour types.
\end{remark}

% ============================================================
\section{4.\enspace Temperature and urgency}
% ============================================================

For a game $G=\{G^L\mid G^R\}$ with $G^L < G^R$, the \emph{temperature}
and \emph{mean value} satisfy
\[
  \temp{G} = \frac{\mean{G^L}-\mean{G^R}}{2}, \qquad
  \mean{G} = \frac{\mean{G^L}+\mean{G^R}}{2}.
\]
In any compound position $\sum_i G_i$, optimal play moves first in the
component with the \emph{largest} temperature.

\begin{proposition}\label{prop:temp}
For the symmetric switch $\pm n = \{n\mid{-n}\}$, we have
$\temp{\pm n} = n$ and $\mean{\pm n}=0$.
The temperature of a Type-III symmetric branch equals the depth
of its branches.
\end{proposition}

\noindent\textbf{Strategic consequence.}
In the compound game $(\pm 2)+(\pm 1)+{*}$, optimal play demands
moving in $\pm 2$ first (hottest, $t=2$), then $\pm 1$ ($t=1$),
treating $*$ last ($t=0$).
Deviating---e.g., playing in $*$ first---grants the opponent a swing of
$2t(\pm 2)=4$ half-moves, which is typically game-deciding.

% ============================================================
\section{5.\enspace Incomparability and the infinitesimal zoo}
% ============================================================

Type-IV positions expose the most striking feature of CGT: the game
order is a \emph{partial} order, not a total one.

\begin{theorem}[Incomparability of $\uparrow$ and $*$]\label{thm:incomp}
$\uparrow\parallel *$, i.e.\ $\uparrow$ and $*$ are incomparable under
the standard CGT ordering.
\end{theorem}

\begin{proof}[Sketch]
Compute $\uparrow - * = \uparrow + * = \{0\mid *\}+*=\{**\mid *\}$.
If Left moves first she plays to $**=0$, winning; if Right moves first
he plays to $*$, leaving $*+*=0$, winning.  Thus $\uparrow+*$ is
\emph{fuzzy} (outcome class $\mathcal{N}$), so $\uparrow\neq *$.
Since $\uparrow>0$ yet $*\parallel 0$, neither $\uparrow>*$ nor
$\uparrow<*$ can hold.
\end{proof}

This means that in positions containing both $\uparrow$ and $*$ summands,
temperature-based heuristics fail entirely: no real-valued urgency ranking
resolves the play correctly, and the full surreal-game machinery is
required.

% ============================================================
\section{6.\enspace Open problems}
% ============================================================

\begin{enumerate}[leftmargin=*,topsep=2pt,itemsep=1pt,label=(\roman*)]
\item \textbf{Complexity.} Is determining the outcome class of a
  Partisan Tree Ascent position in $\mathbf{P}$ or $\mathbf{PSPACE}$-hard?
\item \textbf{Realizability.} Characterise which CGT values arise as
  $\val{v}$ for some edge-labelled binary tree.
\item \textbf{Generalisation.} Replace binary trees with $k$-ary trees
  and $k$ player-colour edge labels.  Do temperature strategies remain
  optimal, or does the richer structure force new phenomena?
\end{enumerate}

\end{document}